
% SGA 3，报告 VII_B；英文重建本第一轮的中文译本。
% 范围： 节 3.3 至 Cartier 定理及其证明结束,
% 依据文献局部页 69--71; 止于下节之前： 节 3.3.1.
% 合并阅读本页： 582--584. 重版印刷页： 571--573.
% 范围内原版边码： 535--538.
% 范围内编者注： 125--128.
% 依据： Polo--Gille Exp7B-13oct24.pdf,
% SHA-256 D303690E0195654BBE7C3DAA5D89EA23E24CD709A3575217234AC19AE8836694.
% 参照了 Jacob Reinhold 的 CC BY 4.0 Markdown，作为对照来源。
% 本组件无图；所有数学内容均为原生 LaTeX。
\subsection*{3.3}
\phantomsection
\label{sga3:VIIB:3.3}
若 \(G\) 是仿射代数为 \(A\) 的形式 \(k\)-群，回顾，\(I_A\) 表示 \(A\) 的增广理想，且
\[
  \omega_{G/k}
  =
  I_A/I_A^2
  \simeq
  I_A\widehat\otimes_Ak
\]
是伪紧 \(k\)-模。
\medskip
\noindent\textbf{定理 3.3 \textup{（Cartier）}。}
\phantomsection
\label{sga3:VIIB:3.3:cartier-theorem}
{\itshape 设 \(k\) 是包含 \(\mathbb Q\) 的伪紧环，\footnote{\label{sga3:VIIB:nde:125}编者删去了 \(k\) 为局部环的假设；证明可归约到这一情形。}\(G\) 是形式 \(k\)-群。下列断言等价：
\begin{enumerate}
\item[\textup{\SGAInlineRefTarget{sga3:VIIB:3.3:cartier-theorem:item:i}(i)}] 存在在 \(\mathcal O_k\) 上平坦的 \(\mathcal O_k\)-Lie 代数 \(\mathbf L\)，使得 \(G\) 同构于 \(\mathscr G(\mathbf L)\)。此时由 \SGARefLink{sga3:VIIB:3.2:prop}{3.2} 有 \(\mathbf L=\operatorname{Lie}(G)\)。
\item[\textup{\SGAInlineRefTarget{sga3:VIIB:3.3:cartier-theorem:item:ii}(ii)}] 存在投射伪紧 \(k\)-模 \(\omega\)，使得 \(G\) 的底形式簇同构于
\[
  \mathbf V_k^{f,0}(\omega)
\]
\textup{（参见\ \SGARefLink{sga3:VIIB:1.2.5}{1.2.5}）}，后者的仿射代数为 \(k[[\omega]]\)。此时 \(\omega\simeq\omega_{G/k}\)。
\item[\textup{\SGAInlineRefTarget{sga3:VIIB:3.3:cartier-theorem:item:iii}(iii)}] 群 \(G\) 无穷小，且 \(\omega_{G/k}\) 是投射伪紧 \(k\)-模。
\item[\textup{\SGAInlineRefTarget{sga3:VIIB:3.3:cartier-theorem:item:iv}(iv)}] 群 \(G\) 无穷小且在 \(k\) 上拓扑平坦。
\end{enumerate}}
\noindent\textit{\textup{\SGARefLink{sga3:vii:paragraph:5-1:equation:i}{(i)}}
\(\boldsymbol{\Rightarrow}\) \textup{\SGARefLink{sga3:vii:paragraph:5-1:equation:ii}{(ii)}} 的证明。}
令
\[
  \omega=\Gamma^*(\mathbf L)
\]
为 \(\mathbf L\) 的对偶伪紧 \(k\)-模\textup{（参见\ \SGARefLink{sga3:VIIB:1.2.3.D}{1.2.3.D}）}。对每个 \(C\in\operatorname{Alf}/k\)，必须构造关于 \(C\) 函子性的同构
\[
  \mathbf V_k^{f,0}(\omega)(C)
  \xrightarrow{\ \sim\ }
  \mathscr G(\mathbf L)(C).
\]
由 \SGARefLink{sga3:VIIB:1.2.5}{1.2.5}，左边可与
\[
  \operatorname{Hom}_c\bigl(\omega,\mathfrak r(C)\bigr),
\]
等同，其中 \(\mathfrak r(C)\) 是 \(C\) 的根。这个集合自然同构于连续 \(C\)-线性映射的集合
\[
  \operatorname{Hom}_c\bigl(
    \omega\widehat\otimes_kC,\mathfrak r(C)
  \bigr)
\]
由于 \(\omega\widehat\otimes_kC\) 是投射伪紧 \(C\)-模，典范映射
\[
  \bigl(\omega\widehat\otimes_kC\bigr)^\dagger
  \otimes_C\mathfrak r(C)
  \longrightarrow
  \operatorname{Hom}_c\bigl(
    \omega\widehat\otimes_kC,\mathfrak r(C)
  \bigr)
\]
是双射\textup{（参见\ \SGARefLink{sga3:VIIB:0.3.6.A}{0.3.6.A}）}。由 \SGARefLink{sga3:VIIB:1.2.3.E}{1.2.3.E}，
\[
  \mathbf L(C)
  \simeq
  \mathbf V_k^f\bigl(\Gamma^*(\mathbf L)\bigr)(C)
  =
  \bigl(\omega\widehat\otimes_kC\bigr)^\dagger.
\]
因此 \(\mathbf V_k^{f,0}(\omega)(C)\) 典范同构于
\[
  \mathbf L(C)\otimes_C\mathfrak r(C),
\]
而\SGARefLink{sga3:VIIB:3.2:prop}{命题 3.2} 将它与 \(\mathscr G(\mathbf L)(C)\) 典范等同。这就证明了\textup{\SGARefLink{sga3:vii:paragraph:5-1:equation:i}{(i)}} \(\Rightarrow\) \textup{\SGARefLink{sga3:VIIB:3.2:prop:item:ii}{(ii)}}。
\medskip
\noindent\textit{\textup{\SGARefLink{sga3:vii:paragraph:5-1:equation:ii}{(ii)}}
\(\boldsymbol{\Rightarrow}\) \textup{\SGARefLink{sga3:VIIB:3.3:cartier-theorem:item:iii}{(iii)}} 的证明。}
设 \(\omega\) 是 \(\operatorname{PC}(k)\) 的投射对象，且
\[
  h:k[[\omega]]\xrightarrow{\ \sim\ }A
\]
是到 \(G\) 的仿射代数的同构。复合
\[
  \varepsilon_A\circ h:k[[\omega]]\longrightarrow k
\]
由其限制 \(\lambda:\omega\to\mathfrak r(k)\) 决定。由 \(k[[\omega]]\) 的泛性质，从 \(\omega\) 到 \(k[[\omega]]\) 的根的映射
\[
  x\longmapsto x-\lambda(x)
\]
延拓为 \(k[[\omega]]\) 的自同态 \(\ell_\lambda\)。由于
\[
  \ell_\lambda\circ\ell_{-\lambda}
  =
  \ell_{-\lambda}\circ\ell_\lambda
  =
  \operatorname{id},
\]
映射 \(\ell_\lambda\) 是自同构。用 \(h\circ\ell_\lambda\) 替换 \(h\)，可设 \(\varepsilon_A\circ h\) 在 \(k[[\omega]]\) 中由 \(\omega\) 生成的闭理想 \(I\) 上为零。于是 \(h\) 诱导同构
\[
  I/I^2\xrightarrow{\ \sim\ }I_A/I_A^2.
\]
由于 \(I/I^2\simeq\omega\)，得到
\[
  \omega_{G/k}\simeq\omega,
\]
故 \(\omega_{G/k}\) 是投射的。形式簇 \(\mathbf V_k^{f,0}(\omega)\)，从而 \(G\)，都是无穷小的。
\medskip
\noindent\textit{\textup{\SGARefLink{sga3:VIIB:3.3:cartier-theorem:item:iii}{(iii)}}
\(\boldsymbol{\Rightarrow}\) \textup{\SGARefLink{sga3:vii:paragraph:5-1:equation:i}{(i)}} 的证明。}
设 \(G\) 无穷小，\(\omega_{G/k}\) 是投射的。令 \(\mathbf L\) 为 \(G\) 的 \(\mathcal O_k\)-Lie 代数。由 \SGARefLink{sga3:VIIB:2.6.1}{2.6.1}，其底 \(\mathcal O_k\)-模为
\[
  \mathbf V_k^f(\omega_{G/k}).
\]
因此，\(\mathbf L\) 在 \(\mathcal O_k\) 上平坦，且\SGARefLink{sga3:VIIB:1.2.3.E}{命题 1.2.3.E} 给出
\[
  \Gamma^*(\mathbf L)\simeq\omega_{G/k}.
\]
所以，\textup{\SGARefLink{sga3:vii:paragraph:5-1:equation:i}{(i)}} \(\Rightarrow\) \textup{\SGARefLink{sga3:vii:paragraph:5-1:equation:ii}{(ii)}} 的证明将 \(\mathscr G(\mathbf L)\) 的仿射代数与 \(k[[\omega_{G/k}]]\) 等同。由 \SGARefLink{sga3:VIIB:2.6.3:prop}{2.6.3}，\(\mathbf L\) 的恒等态射典范地对应于一个态射
\[
  \mathscr G(\mathbf L)\longrightarrow G,
\]
从而对应于连续 \(k\)-代数态射
\[
  \varphi:A\longrightarrow k[[\omega_{G/k}]].
\]
令 \(I\) 为 \(\omega_{G/k}\) 生成的闭理想。用幂 \(I^n\) 的闭包滤过 \(k[[\omega_{G/k}]]\)，用幂 \(I_A^n\) 的闭包滤过 \(A\)。由定义，\(\varphi\) 在
\[
  \omega_{G/k}=I_A/I_A^2=I/I^2.
\]
上诱导恒等映射。
由于 \(\omega_{G/k}\) 是投射的，典范投影
\[
  I_A\longrightarrow\omega_{G/k}
\]
有一个截面 \(\tau\)。由 \(k[[\omega_{G/k}]]\) 的泛性质，这个截面定义了连续代数态射
\[
  \psi:k[[\omega_{G/k}]]\longrightarrow A.
\]
复合 \(\varphi\circ\psi\) 在 \(\omega_{G/k}\) 上诱导恒等映射，从而在 \(k[[\omega_{G/k}]]\) 的相伴分次代数上也诱导恒等映射。由\textup{[CA]，\S V.7，引理 1}，\(\varphi\circ\psi\) 是同构。\footnote{\label{sga3:VIIB:nde:126}比较下文 \SGARefLink{sga3:VIIB:5.1.5}{5.1.5}，以及\textup{[BAC]，III，\S \SGARefLink{sga3:VIIB:2.8}{2.8}，定理 1 及其推论}。以下原论证的细节由编者补充。}
此外，\(\psi\) 诱导同构
\[
  I/I^2\xrightarrow{\ \sim\ }I_A/I_A^2,
\]
从而在相伴分次代数上诱导满射。由于 \(G\) 是纯不可分的（radicial），\(I_A\) 包含于 \(A\) 的根中，故
\[
  \bigcap_n I_A^n=0.
\]
所引引理因此表明 \(\psi\) 是满射。由于 \(\varphi\circ\psi\) 是同构而 \(\psi\) 是满射，\(\psi\) 和 \(\varphi\) 都是同构。这证明了\textup{\SGARefLink{sga3:VIIB:3.3:cartier-theorem:item:iii}{(iii)}} \(\Rightarrow\) \textup{\SGARefLink{sga3:VIIB:3.3:cartier-theorem:item:i}{(i)}}。
\medskip
最后，\textup{\SGARefLink{sga3:vii:paragraph:5-1:equation:i}{(i)}} 或\textup{\SGARefLink{sga3:vii:paragraph:5-1:equation:ii}{(ii)}} 显然蕴含\textup{\SGARefLink{sga3:VIIB:3.3:cartier-theorem:item:iv}{(iv)}}，所以只须再证\textup{\SGARefLink{sga3:VIIB:3.3:cartier-theorem:item:iv}{(iv)}} \(\Rightarrow\) \textup{\SGARefLink{sga3:VIIB:3.3:cartier-theorem:item:ii}{(ii)}}。可设 \(k\) 是局部环，剩余域为 \(k_0\)。置
\[
  G_0=G\widehat\otimes_kk_0,
  \qquad
  \omega=\omega_{G/k},
  \qquad
  \omega_0=\omega\widehat\otimes_kk_0.
\]
\footnote{\label{sga3:VIIB:nde:127}以下原论证的细节由编者补充。}
伪紧 \(k_0\)-模 \(\omega_0\) 有伪基 \((y_\lambda)_{\lambda\in\Lambda}\)。令
\[
  \omega'=\prod_{\lambda\in\Lambda}k_\lambda
\]
为以该集合为指标集的拓扑自由 \(k\)-模，并将每个 \(y_\lambda\) 提升为元素 \(x_\lambda\in\omega\)。这给出连续 \(k\)-线性映射
\[
  f:\omega'\longrightarrow\omega
\]
使得
\[
  f_0=f\widehat\otimes_kk_0
\]
可逆。\footnote{\label{sga3:VIIB:nde:128}并不能预先知道 \(\omega\) 是投射伪紧 \(k\)-模；这将从论证中得到。比较 VIIA~\SGARefLink{sga3:VIIA:7.4}{7.4} 中\textup{\SGARefLink{sga3:VIIB:3.3:cartier-theorem:item:iv}{(iv)}} \(\Rightarrow\) \textup{\SGARefLink{sga3:VIIB:3.3:cartier-theorem:item:iii}{(iii)}} 的证明。}
由于 \(\omega'\) 是投射的，\(f\) 提升为连续 \(k\)-线性映射
\[
  g:\omega'\longrightarrow I_A
\]
满足 \(\pi\circ g=f\)，其中 \(\pi:I_A\to\omega\) 为投影。\(k[[\omega']]\) 的泛性质给出拓扑代数态射
\[
  \phi:k[[\omega']]\longrightarrow A.
\]
现在
\[
  \omega'\widehat\otimes_kk_0
  \simeq
  \omega\widehat\otimes_kk_0
  =
  \omega_{G_0/k_0},
\]
因而
\[
  k[[\omega']]\widehat\otimes_kk_0
  \simeq
  k_0[[\omega_{G_0/k_0}]].
\]
对 \(k_0,G_0\) 满足\textup{\SGARefLink{sga3:VIIB:3.3:cartier-theorem:item:iii}{(iii)}}，而\textup{\SGARefLink{sga3:VIIB:3.3:cartier-theorem:item:iii}{(iii)}} \(\Rightarrow\) \textup{\SGARefLink{sga3:vii:paragraph:5-1:equation:i}{(i)}} 的证明表明
\[
  \phi_0=\phi\widehat\otimes_kk_0
\]
可逆。\(k[[\omega']]\) 与 \(A\) 都是投射伪紧 \(k\)-模，所以 \SGARefLink{sga3:VIIB:0.3.4}{0.3.4} 蕴含 \(\phi\) 本身可逆。特别地，若 \(I\) 是 \(k[[\omega']]\) 的增广理想，则 \(\phi\) 诱导同构
\[
  \omega'=I/I^2
  \xrightarrow{\ \sim\ }
  I_A/I_A^2=\omega.
\]
证明完毕。
